\section{仿真结果与输出耦合}

\subsection{端接条件对等效关系的影响}

表~\ref{tab:branches} 汇总三组支路构造。各组等式与不等式关系均在扫频前通过代数运算确定。

\begin{table}[htbp]
\centering
\caption{归一化元件值与支路系数。}
\label{tab:branches}
\begin{tabular}{lrrrrrrrr}
\toprule
支路 & $r_1$ & $r_2$ & $q_1$ & $q_2$ & $a$ & $b$ & $c$ & $d$ \\
\midrule
P & 1 & 1 & 1 & 1 & 3 & 1 & 2 & 1 \\
Q & 1 & 4 & $1/2$ & $1/2$ & 3 & 1 & 5 & 2 \\
T & $18/19$ & $20/19$ & $361/180$ & $1/2$ & 2.9 & 1 & 2 & 2 \\
\bottomrule
\end{tabular}
\end{table}

图~\ref{fig:termination-transfer} 比较 P/Q 的开路和带载响应。图~\ref{fig:output-impedance} 给出区分这两条开路等效支路的输出阻抗。图~\ref{fig:load-specific} 比较两种负载下 T 相对于 P 的响应。



仿真复现了 P/Q 的开路共模抵消，以及 P/T 在 \(R_L=10\,\mathrm{k\Omega}\) 下的共模抵消。将 \(R_L\) 改为 5 k\(\Omega\) 后，P/T 响应不再相等。在全部 37 个算例中，LTspice 输出与独立节点方程解之间的复节点电压绝对差最大值为 \(1.14\times10^{-15}\) V。

\begin{figure}[htbp]
\centering
\includegraphics[width=0.67\linewidth]{figures/Fig_03_termination_transfer.pdf}
\caption{P、Q 在开路和 10~k$\Omega$ 对地负载下的幅频响应。实线为解析传递函数，空心圆为已有 LTspice 仿真数据。两条开路响应重合。归一化基准为 $R_*=1$~k$\Omega$、$C_*=10$~nF。}
\label{fig:termination-transfer}
\end{figure}

\begin{figure}[htbp]
\centering
\includegraphics[width=0.67\linewidth]{figures/Fig_04_output_impedance.pdf}
\caption{理想输入电压源置零后，P、Q、T 的开路输出阻抗幅值。曲线由四系数表达式计算。P 与 Q 的开路电压传递函数相同，输出阻抗不同。}
\label{fig:output-impedance}
\end{figure}

\begin{figure}[htbp]
\centering
\includegraphics[width=0.67\linewidth]{figures/Fig_05_load_specific_match.pdf}
\caption{两种共用对地负载下 T 相对于 P 的幅值响应。纵轴为 $20\log_{10}|H_T/H_P|$，零值表示幅值相等。实线为解析结果，空心圆为已有 LTspice 仿真数据。10~k$\Omega$ 下完整传递函数的相等性由正文代数推导确定。}
\label{fig:load-specific}
\end{figure}


\subsection{跨接电容同时作用于两种差分输出增益}

保持理想电压源和指定并联负载不变，在两个输出节点之间接入电容 \(C_d\)。以 \(Z'_P\) 和 \(Z'_N\) 分别表示各支路输出阻抗与本侧负载并联后的阻抗，以 \(H_P\) 和 \(H_N\) 表示接入 \(C_d\) 前的带载传递函数。联立两个输出节点方程可得

\begin{equation}\label{eq:12}
\begin{aligned}
G_{DC}&=\frac{H_P-H_N}{D_d},\\
G_{DD}&=\frac{(H_P+H_N)/2}{D_d},\\
D_d&=1+sC_d(Z'_P+Z'_N).
\end{aligned}
\end{equation}

在 \(G_{DD}/G_{DC}\) 有定义的频点，共同因子 \(D_d\) 从比值中约去。跨接电容以同一因子改变两种差分输出增益，因而保持 \(G_{DD}/G_{DC}\) 不变。若原带载支路传递函数相等，接入跨接电容后仍有 \(G_{DC}\) = 0。

对于 \(R_L=10\,\mathrm{k\Omega}\)、\(C_d=820\,\mathrm{pF}\) 的 P/Q 网络，在 145 个频点上，相对于未接跨接电容时复增益比的最大相对差为 \(5.14\times10^{-12}\)。该结果与式~\eqref{eq:12} 中共同因子的约去一致。


图~\ref{fig:bridge-gains} 展示两个独立增益的变化，图~\ref{fig:bridge-cmrr} 比较其共模抑制比。

\begin{figure}[htbp]
\centering
\includegraphics[width=0.67\linewidth]{figures/Fig_06_bridge_gains.pdf}
\caption{P/Q 支路对在有、无跨接电容时的解析差分输出增益。DM 表示单位差模输入；CM$\rightarrow$DM 表示单位共模输入产生的差分输出。空心标记用于识别 820~pF 曲线，其数值同样来自解析计算。两侧对地负载均为 10~k$\Omega$。}
\label{fig:bridge-gains}
\end{figure}


\begin{figure}[htbp]
\centering
\includegraphics[width=0.67\linewidth]{figures/Fig_07_bridge_cmrr.pdf}
\caption{负载条件下 P/Q 支路对的共模抑制比。$C_d=0$ 和 820~pF 的解析曲线重合。空心圆由 $C_d=820$~pF 时已有共模和差模 LTspice 输出差值计算。两侧对地负载均为 10~k$\Omega$；CMRR 为 $20\log_{10}|G_{DD}/G_{DC}|$。图~\ref{fig:bridge-gains} 同时给出绝对增益，图~\ref{fig:terminal-sweep} 展示独立支路的端接扫描。}
\label{fig:bridge-cmrr}
\end{figure}


\FloatBarrier
